\section{The Strongly Convex Case}
\label{sec:strongly-convex}

In this section, we prove that in the case of strong convexity, our algorithm can be improved by regularizing the loss of online learner.
If $ f$ is strongly convex, then we can prove a similar result to Theorem \ref{lem:basic-lemma} (the proof is in Appendix \ref{app:strongly-convex}).

\begin{lemma}
\label{lem:DP-OTB:strongly-convex-regularized-loss}
Suppose $f$ is $\mu$-strongly convex w.r.t. $\|\cdot\|$. If we replace $\ell_t(w)=\langle g_t+\gamma_t,w\rangle$ with $\bar \ell_t(w)=\ell_t(w)+\frac{\weight_t\mu}{4}\|w-x_t\|^2$ in Algorithm \ref{alg:DP-OTB:main}, and denote the associated regret by $\overline \regret_T$, then
\begin{align*}
    \weight_{1:T} ( f(x_T) -  f(x^*))
    &\leq \overline\regret _T(x^*) +\sum_{t=1}^T \langle \weight_t\nabla f(x_t)-g_t-\gamma_t, w_t-x^*\rangle - \frac{\weight_t\mu}{8}\|w_t-x^*\|^2 \\
    &\leq \overline\regret_T(x^*) + \sum_{t=1}^T \frac{2\|\weight_t\nabla f(x_t)-g_t-\gamma_t\|_*^2}{\weight_t \mu}.
\end{align*}
\end{lemma}

% First, recall that $x_t$ are determined by $\{g_1+\gamma_1,\ldots,g_{t-1}+\gamma_{t-1}\}$, and we proved in Theorem \ref{thm:main-privacy}

Compared to Lemma \ref{lem:basic-lemma} and Equation \ref{eq:basic-decomposition-2}, in the strongly convex case, there is an additional term $-\frac{\weight_t\mu}{8}\|w_t-x^*\|^2$, which allows the improved convergence rate:

\begin{theorem}
\label{thm:strongly-convex}
Suppose Assumption \ref{as:sc-norm} - \ref{as:sigma-H} hold, and $\calD$ is a $(V,\alpha)$-RDP distribution. Also suppose $ f$ is $\mu$-strongly convex. Set $\weight_t=t^k$ and $\sigma_t^2$ as defined in \eqref{eq:RDP-noise}. Then $\E[ f(x_T) -  f(x^*)]$ is bounded by:
\begin{align*}
    % \E[ f(x_T)- f(x^*)] &\leq 
    \frac{(k+1)\E[\overline\regret_T(x^*)]}{T^{k+1}} 
    + \frac{16(k+1)^3}{\lambda \mu} \left( \frac{(\sigma_G+D\sigma_H)^2}{T} + \frac{2V(G+DH)^2\log_2^2(2T)}{\rho^2T^2} \right).
\end{align*}
Moreover, for all $\bar g_t\in \partial \bar \ell_t(w_t)$,
\begin{align*}
    \E[\|\bar g_t\|_*^2] \leq t^{2k}\left( 4G^2 + \mu^2D^2 + \frac{16(k+1)^2}{\lambda}\left(\frac{(\sigma_G+D\sigma_H)^2}{t} + \frac{2V(G+DH)^2\log_2^2(2T)}{\rho^2t^2}\right) \right).
\end{align*}
\end{theorem}

\begin{remark}
As an example, again consider the case $\calD=\calN(0,I)$ with $L_2$ norm. Online subgradient descent (OSD) with appropriate learning rate on $\mu_t$-strongly convex losses $\ell_t$ achieves:
% the following regret:
\begin{align*}
    \regret_T(u) \leq \sum_{t=1}^T \frac{\|g_t\|_2^2}{2\sum_{i=1}^t\mu_i},
\end{align*}
where $g_t\in \partial\ell_t(w_t)$. In our case, $\|\cdot\|_2^2$ is $2$-strongly convex and the regularized loss $\bar \ell_t$ is $\frac{\weight_t\mu}{2}$-strongly convex, i.e. $\mu_t=\mu t^k/2$ and $\sum_{i=1}^t \mu_i = O(\mu t^{k+1})$. Therefore, by Theorem \ref{thm:strongly-convex},
\begin{align*}
    \E[\overline\regret_T(x^*)] 
    % &\lesssim \sum_{t=1}^T \frac{t^{k-1}}{\mu}\left(G^2+\mu^2D^2+ \frac{(\sigma_G+D\sigma_H)^2}{t}+\frac{\alpha(G+DH)^2\log_2^{3/2}T}{\rho^2t^2}\right) \\
    % &\leq \frac{T^k}{\mu}\left(G^2 + \mu^2 D^2 + \frac{(\sigma_G+D\sigma_H)^2}{T}+\frac{\alpha(G+DH)^2\log_2^{3/2}T}{\rho^2T^2} \right).
    &\leq O\left( \frac{T^k}{\mu}\left( G^2+\mu^2D^2 + \frac{(\sigma_G+D\sigma_H)^2}{T} + \frac{d(G+DH)^2\log^2T}{\rho^2T^2}  \right) \right).
\end{align*}
Consequently,
\begin{align*}
    \E[f(x_T)- f(x^*)]
    \leq O\left(  \frac{(G+\mu D+D\sigma_H)^2}{\mu T} + \frac{d(G+DH)^2\log^2T}{\mu\rho^2T^2} \right).
\end{align*}
This again matches the optimal private convergence rates.
\end{remark}

\begin{proof}[Proof of Theorem \ref{thm:strongly-convex}]
We have already bounded $\E[\|\weight_t\nabla f(x_t)-g_t\|_*^2]$ and $\E[\|\gamma_t\|_*^2]$ in \eqref{eq:variance-norm} and \eqref{eq:gamma-bound} respectively, so
\begin{align*}
    \E[\|\weight_t\nabla f(x_t)-g_t-\gamma_t\|_*^2]
    &\leq 2\E[\|\weight_t\nabla f(x_t)-g_t\|_*^2+\|\gamma_t\|_*^2] \\
    % &\leq 8(k+1)^2(\sigma_G^2+D^2\sigma_H^2)t^{2k-1}/\lambda, \label{eq:variance-norm} + \frac{4K\log_2(2t)}{\lambda}\left(\frac{2(k+1)(G+DH)t^{k-1}\log_2(2T)}{\rho}\right)^2.
    &\leq \frac{8(k+1)^2t^{2k-1}}{\lambda}\left((\sigma_G+D\sigma_H)^2+ \frac{2V(G+DH)^2\log_2^2(2T)}{\rho^2 t}\right).
\end{align*}
Upon substituting this into Lemma \ref{lem:strongly-convex-regularized-loss} and replace $\weight_t=t^k$, we get:
\begin{align*}
    & \weight_{1:T} \E[ f(x_T)- f(x^*)] \\
    \leq & \E[\overline\regret_T(x^*)] + \sum_{t=1}^T \frac{16(k+1)^2t^{k-1}}{\lambda\mu}\left((\sigma_G+D\sigma_H)^2+ \frac{2V(G+DH)^2\log_2^2(2T)}{\rho^2 t}\right) \\
    % \intertext{Recall $\weight_t=t^k$ and $\weight_{1:T}\geq T^{k+1}/(k+1)$.}
    \leq & \E[\overline\regret_T(x^*)] + \frac{16(k+1)^2}{\lambda \mu} \left( (\sigma_G+D\sigma_H)^2T^k + \frac{2V(G+DH)^2\log_2^2(2T)T^{k-1}}{\rho^2} \right).
\end{align*}
Dividing both sides by $\weight_{1:T} \geq T^{k+1}/(k+1)$ proves the first part of the theorem.

For the second part, recall that $\bar \ell_t(w)=\langle g_t+\gamma_t,w\rangle + \frac{\weight_t\mu}{4}\|w-x_t\|^2$. Therefore, for all $\bar g_t\in\partial \bar\ell_t(w)$, $\bar g_t = g_t+\gamma_t + \frac{\beta_t\mu}{4}v$, where $v\in \partial\|w_t-x_t\|^2$.
% Since the linear part is differentiable,
% \begin{align*}
%     \partial\bar{\ell}_t(w) = g_t+\gamma_t + \frac{\weight_t\mu}{4}\partial\|w-x_t\|^2. 
% \end{align*}
% By chain rule of sub-differentials (Corollary 16.72 \cite{bauschke2011convex}), let $\phi(r)=r^2$ and $f(w)=\|w-x_t\|$, then $\|w-x_t\|^2 = \phi\circ f(w)$ and
% \begin{align*}
%     \partial \|w-x_t\|^2 
%     &= \{\alpha u : \alpha\in \partial \phi(f(w)), u\in\partial f(w) \} \\
%     &= \{2\|w-x_t\| u : u\in\partial \|w-x_t\|\}.
% \end{align*}
% By assumption, $\|w-x_t\|\leq D$. By triangular inequality, $\|x\|-\|y\| \leq \|x-y\|$, and thus $\|\cdot\|$ is $1$-Lipschitz. Equivalently, $\|u\|_*\leq 1$ for all $u\in\partial \|w-x_t\|$. Therefore, for all $v\in \partial\|w-x_t\|^2$, $\|v\|_*\leq 2D$.
We follow the same argument in Theorem \ref{thm:DP-OTB-general-norm}: 
\begin{align*}
    \E[\|\bar g_t\|_*^2]
    &\leq 4\E[\|\weight_t\nabla f(x_t)\|_*^2 + \|\weight_t\nabla f(x_t)-g_t\|_*^2 + \|\gamma_t\|_*^2 + \|\tfrac{\weight_t\mu}{4} v\|_*^2] \\
    \intertext{Here $v\in\partial\|w_t-x_t\|^2$. We bound the first term by Lipschitz, the second by \eqref{eq:variance-norm}, and the third by \eqref{eq:gamma-bound}. Moreover, by chain rule (Proposition \ref{prop:subdiff-chain-rule}), we can show that $\|v\|_* \leq 2D$, so:}
    &\leq 4t^{2k}G^2 + \frac{16(k+1)^2(\sigma_G+D\sigma_H)^2t^{2k-1}}{\lambda} \\
    &\quad \ + \frac{32V(k+1)^2(G+DH)^2t^{2k-2}\log_2^2(2T)}{\lambda\rho^2} + \mu^2D^2t^{2k} \\
    &\leq t^{2k}\left( 4G^2 + \mu^2D^2 + \frac{16(k+1)^2}{\lambda}\left(\frac{(\sigma_G+D\sigma_H)^2}{t} + \frac{2V(G+DH)^2\log_2^2(2T)}{\rho^2t^2}\right) \right).
\end{align*}
\end{proof}

